\section{Conclusiones}

\begin{frame}{Resumen del Pipeline 3DGS}
  \begin{tikzpicture}[node distance=0.6cm, every node/.style={font=\small}]
    \node[draw=MidnightBlue, fill=MidnightBlue!15, rounded corners, minimum width=2.8cm, minimum height=0.75cm] (sfm)    at (0,0)   {SfM $\to$ nube de puntos};
    \node[draw=MidnightBlue, fill=MidnightBlue!20, rounded corners, minimum width=2.8cm, minimum height=0.75cm] (init)   at (3.5,0)  {Inicialización Gaussianas};
    \node[draw=MidnightBlue, fill=MidnightBlue!25, rounded corners, minimum width=2.8cm, minimum height=0.75cm] (raster) at (7,0)    {Proyección 2D + Raster.};
    \node[draw=MidnightBlue, fill=MidnightBlue!30, rounded corners, minimum width=2.8cm, minimum height=0.75cm] (loss)   at (10.5,0) {$\mathcal{L}_1 + \mathcal{L}_\text{SSIM}$};
    \node[draw=ForestGreen,  fill=ForestGreen!15,  rounded corners, minimum width=2.8cm, minimum height=0.75cm] (adam)   at (7,-1.5) {Adam + densificación};
    \node[draw=accentorange, fill=accentorange!15, rounded corners, minimum width=2.8cm, minimum height=0.75cm] (out)    at (3.5,-1.5){Escena final $>$90 fps};

    \draw[->, thick] (sfm)    -- (init);
    \draw[->, thick] (init)   -- (raster);
    \draw[->, thick] (raster) -- (loss);
    \draw[->, thick] (loss)   -- (adam);
    \draw[->, thick] (adam)   -- (out);
  \end{tikzpicture}
\end{frame}

\begin{frame}{Contribuciones Principales}
  \begin{columns}
    \begin{column}{0.5\textwidth}
      \textbf{Matemáticas clave:}
      \begin{itemize}
        \item Factorización $\boldsymbol{\Sigma} = \mathbf{RSS}^\top\mathbf{R}^\top$ con cuaterniones
        \item Armónicos Esféricos $L=3$ para color dependiente de vista
        \item Alpha blending discreto equivalente a NeRF
      \end{itemize}
    \end{column}
    \begin{column}{0.5\textwidth}
      \textbf{Sistema de entrenamiento:}
      \begin{itemize}
        \item Pérdida combinada $\mathcal{L}_1 + \mathcal{L}_{\text{D-SSIM}}$
        \item Control adaptativo de densidad (clonar/dividir/podar)
        \item Reset periódico de opacidades
      \end{itemize}
    \end{column}
  \end{columns}
  \vspace{1em}
  \begin{exampleblock}{Resultado}
    3DGS logra calidad comparable a NeRF con renderizado en \textbf{tiempo real} y entrenamiento de \textbf{20--40 minutos}.
  \end{exampleblock}
\end{frame}

\begin{frame}{Trabajo Futuro y Limitaciones}
  \begin{columns}
    \begin{column}{0.5\textwidth}
      \textbf{Limitaciones actuales:}
      \begin{itemize}
        \item Alto coste de memoria ($\sim$1 GB/escena)
        \item Artefactos en regiones sin cobertura de cámara
        \item Las secciones de proyección 2D y rasterización requieren mayor desarrollo
      \end{itemize}
    \end{column}
    \begin{column}{0.5\textwidth}
      \textbf{Líneas de investigación:}
      \begin{itemize}
        \item Compresión de primitivas Gaussianas
        \item Escenas dinámicas (4DGS)
        \item Integración con LiDAR para mejor inicialización
        \item Edición semántica de escenas 3DGS
      \end{itemize}
    \end{column}
  \end{columns}
\end{frame}

\begin{frame}[standout]
  \centering
  \LARGE \textbf{¿Preguntas?}\\[1em]
  \normalsize
  Jaime Hernández · Yauheni Yatsevich\\
  Universidad de Alicante, Marzo 2026
\end{frame}
